\subsection{Definición de los autómatas celulares}
\label{section:def}

Los \emph{autómatas celulares} (AC) son sistemas dinámicos, discretos en el espacio y en el tiempo, caracterizados por la interacción local y una forma de evolución inherentemente paralela \cite{ilachinski2001cellular}. Estos sistemas están formados por elementos simples y homogéneos llamados \emph{celdas}. Cada celda puede existir (no simultáneamente) en $k$ estados diferentes, donde $k$ es un número finito igual o mayor que $2$. El \emph{enmallado} de celdas puede ser $n$ ($\ge 1$) dimensional. Para que las celdas del enmallado evolucionen, es necesario una \emph{función de transición local} que dicte cómo cada celda altera su estado en base a un conjunto de reglas que tienen en cuenta el estado actual de la celda y el estado actual de las celdas de su \emph{vecindad} \cite{schiff2007cellular}.